\section*{How to read this paper}
\addcontentsline{toc}{section}{How to read this paper}
\markboth{How to read this paper}{How to read this paper}

The proof is short in outline and long in detail, and the details are of three very different
kinds: bookkeeping about subsets, a linear-algebra argument about matchings, and an extremal
choice. They are much easier to absorb one at a time. So:

\begin{description}[leftmargin=1.2em,itemsep=2pt]
\item[Level 1 --- five minutes.] Read \cref{sec:intro}, look at \cref{fig:roadmap}, and read the
  boxed \emph{Idea} passages. You will know what is proved and what the three engines are.
\item[Level 2 --- an hour.] Read \cref{sec:reduction,sec:incompatibility} (which use each other)
  and skip \cref{sec:pruning} on a first pass, taking \cref{lem:pruning} as a black box. Then read
  \cref{sec:pruning} only for the statement and the shape of the argument. \Cref{sec:weighted,sec:cycles}
  are consequences and are not needed for the main theorem.
\item[Level 3 --- a day.] Everything, in order, including the toolkit (\appref{app:toolkit}), which
  proves the five facts taken from outside graph theory, and the verification appendix
  (\appref{app:checks}), which records how each step was checked.
\end{description}

\paragraph{Conventions.}
\begin{itemize}[leftmargin=1.2em,itemsep=1pt]
\item All graphs are finite, and \emph{oriented} means: no loops, no two arcs between the same
  pair of vertices in the same direction, and never both $x\to y$ and $y\to x$.
\item $F(S)$ is the union of the out-neighborhoods of the vertices of $S$; it \emph{never}
  subtracts $S$. This convention is used constantly and is stated again where it matters.
\item $|S|$ is the cardinality of a finite set, $\setminus$ is set difference, and $\sqcup$ is a
  disjoint union (used when the two parts are known to be disjoint).
\item The word \emph{proper} means $\neq$ the whole vertex set; \emph{nonempty proper} subsets are
  the only subsets that ever appear in the subset inequalities.
\end{itemize}

\paragraph{The four kinds of box.} Each has a fixed job, so you can decide at a glance whether to
read it.
\begin{itemize}[leftmargin=1.2em,itemsep=1pt]
\item A \emph{What this section does} box opens every section: assumptions in, result out, and why
  the section exists.
\item An \emph{Idea} box says why a definition or a strategy is the right one. It is motivation,
  not part of any proof.
\item A \emph{Watch out} box flags a trap: a hypothesis that is silently used, or a plausible
  statement that is false.
\item A \emph{Where we are} box appears at the junctions of the long argument and says what has
  been established so far.
\end{itemize}
Formal statements are preceded by a grey \emph{In words} line: one sentence without symbols
saying what the statement asserts, so that you can read the statement without parsing it. Long
proofs are preceded by a green \emph{Proof plan} and split into labeled steps.

\paragraph{Notation at a glance.} Every symbol used in the paper appears here, once, with the
place where it is introduced. The dictionary in \appref{app:dictionary} translates back to the
notation of the original preprint.

{\small
\renewcommand{\arraystretch}{1.16}
\noindent\begin{tabularx}{\textwidth}{@{}lX@{}}
\toprule
\textbf{Symbol} & \textbf{Meaning} \\
\midrule
$G$, $D$, $V(G)$, $A(G)$ & an oriented graph, its vertex set, its arc set; $x\to y$ means $(x,y)\in A(G)$ \\
$n$ & the number of vertices (\emph{order}) of the graph under discussion \\
$F(S)$ & $\{y:\ \exists x\in S,\ x\to y\}$; the union of out-neighborhoods, never subtracting $S$ (\cref{sec:intro}) \\
$F^0(S),F^2(S),F^k(S)$ & $F^0(S)=S$, $F^{k+1}(S)=F(F^k(S))$; $F^2=F\circ F$ \\
$N_1^+(v)$, $N_2^+(v)$ & first and \emph{exact} second out-neighborhood: $F(\{v\})$ and $F^2(\{v\})\setminus F(\{v\})$ \\
$d^+(v)$, $d^-(v)$, $\delta^\pm(D)$ & out-degree, in-degree, and their minima over $V(D)$; $d^\pm$ only in \cref{sec:cycles} \\
$S$, $I$, $E$ & a vertex set, and its split $I=F(S)\cap S$, $E=F(S)\setminus S$ (\cref{sec:reduction}) \\
$g(Y)$ & $\abs{F(Y)\setminus(I\cup Y)}$, the counted quantity in \cref{lem:deficit} \\
$Y$, $D_Y$ & a subset of $E$, and $D$ with all arcs $S\to E\setminus Y$ deleted \\
$Q_v$ & $F(\{v\})\cap(E\setminus Y)$: the out-neighbors of $v$ that the deletion above removes \\
$\Lambda$ & $F(Q_v)\setminus(I\cup Y\cup Q_v\cup F(Y))$: second out-neighbors lost by the deletion \\
$m$, $T_m$, $\tau$, $U$ & block size, transitive tournament on $m$ vertices, its internal image operator, and a vertex set of the blowup (\cref{sec:reduction}) \\
$W$ & $F(S)\cup F^2(S)$: the set of blocks that are completely filled (\cref{lem:blowup}) \\
$X$, $X\times X$ & the vertex set of the graph $G$ of \cref{sec:incompatibility}, and the set of ordered pairs \\
$\Phi_1 P$, $\Phi_2 Q$ & coordinate images of a family of pairs: advance the first, resp.\ second, coordinate along an arc \\
$\Delta$ & $\{(v,v):v\in X\}$, the diagonal \\
$P$, $Q$, $M$ & the extremal left and right families and $M=|P|+|Q|$ \\
$P_j$, $Q_i$ & the column $\{p:(p,j)\in P\}$ and the row $\{s:(i,s)\in Q\}$ \\
$u$, $u'$ & the number of uncovered points, before and after pruning \\
$\mathcal L$, $\mathcal R$ & the left and right family of pairs of \cref{lem:pruning}; in \cref{sec:incompatibility} they are the complements of $\Phi_2^2Q$, $\Phi_1^2P$ \\
$\mathcal L'$, $\mathcal R'$ & the subfamilies retained by the pruning lemma \\
$Z$, $H$ & the conflict points $(i,j)$ and the collision pairs $(p,s)$ of \cref{lem:pruning} \\
$q$ & the number of points of $Z$ left uncovered by $\mathcal L'\cup\mathcal R'$ \\
$\Gamma$, $K$, $\mathcal M$ & the augmented bipartite graph, a vertex cover of it, a maximum matching \\
$k$, $\ell$, $t$ & the numbers of original--original, left--$Z$, and $Z$--right edges of $\mathcal M$ \\
$L$, $N$, $A$, $B$ & the four linear maps of the commuting square \eqref{eq:commute} \\
$D^{p,i}$, $T^{p,i}$ & auxiliary matrix for the pair of indices $p\to i$, and its allowed support \\
$J_p$, $S_i$ & $\{j:(p,j)\in\mathcal L\}$ and $\{s:(i,s)\in\mathcal R\}$ \\
$a_{xy}$, $b_{xy}$ & two independent arrays of variables, one per arc $x\to y$ \\
$\overline N(v)$ & the non-neighbors of $v$: $V\setminus(\{v\}\cup N_1^+(v)\cup N_1^-(v))$ (\cref{sec:cycles}) \\
\bottomrule
\end{tabularx}}
